\chapter*{前言}
\addcontentsline{toc}{chapter}{前言}
\markboth{前言}{前言}

本讲义配合一学期的课程《量子科学计算》使用，面向数学、物理与计算机科学专业的高年级
本科生与研究生。其宗旨，是把量子计算的工具组织成一套系统的工具箱，用以攻克科学与工程
中的大规模数值问题——线性方程组、特征值与奇异值问题、矩阵函数、微分方程、哈密顿量
模拟与优化——并在每一步都诚实地交代：量子加速从何而来，又要付出怎样的代价。

阅读本讲义，仅假定读者熟悉线性代数（内积空间、谱定理、奇异值分解）、基本的数值分析
（误差、条件数、迭代法），以及初步的量子力学或量子信息（态、测量、酉演化）；\emph{不}
假定此前修读过量子计算课程——所需的一切将在第一部分从头建立。

全书分为五个部分、共二十讲。第一部分奠定量子计算的基础；第二部分围绕\emph{块编码}与
\emph{量子奇异值变换}这两个统一的概念，发展量子数值线性代数；第三至第五部分则把这些
工具分别投向微分方程与模拟、张量网络，以及优化与量子机器学习。各部分并非彼此独立：
第二部分的块编码与 QSVT 在后续应用中反复出现，读者宜常常回顾。

每一讲都遵循相同的节奏：以一个学习目标方框开篇，经由定义、定理与\emph{完整展开的
例子}——带有显式线路与数值核验——发展其主题，并以习题与延伸阅读收尾。全书的基调刻意
偏重教学：这是一份讲义，宁取清晰的动机与完整的例子，而不尚简省的形式。

本讲义仍在逐讲撰写之中，尚未写就之处已在源文件中标注，粗疏之处尚祈读者包涵。在基础
与量子数值线性代数部分，讲义参考了 Lin Lin 的讲义、Dong An 关于量子数值线性代数的
报告，以及 Childs 的讲义，相关来源在各讲中逐一标明。本讲义为英文版的简体中文翻译，
两版并行维护。欢迎来函指正与建议。
